\begin{abstract}
Knowledge graph (KG) reasoning systems are typically evaluated on link-prediction
accuracy alone, leaving inconsistency between facts both undetected and
undiagnosable. We introduce a sheaf-theoretic framework that treats a KG as a
cellular sheaf over an ontology log (Olog), making logical inconsistency
\emph{quantifiable} as the dimension of the first cohomology group $H^1$.
On top of this, we factor type-safe reasoning into two complementary
mechanisms: a typed attention layer~(B) that gates information flow by Olog
reachability, and a constrained decoder~(D) that restricts the output
distribution to admissible labels. We show empirically and theoretically
that (B) is necessary but not sufficient for trajectory soundness, and that
(B) and (D) target distinct correctness properties: information-flow
correctness inside the model and output validity at the sampling boundary.
On FB15K-237 we obtain MRR \textbf{0.346} and Hits@10 \textbf{0.524}, on
par with KG-BERT, KGT5 ($0.300$), and SimKGC ($0.336$), while adding two
guarantees they lack: (i) a quantitative inconsistency signal---injecting
76 conflicts into a clean graph raises $H^1$ by \textbf{+53}---and (ii) a
typed (B)+(D) pipeline that drives invalid attention weight to zero and
provides a derivation alongside each prediction. WN18RR's tree-structured
ontology yields a sheaf consistency score of \textbf{0.633} versus
\textbf{0.292} on Freebase, correctly reflecting the underlying topology.
Our contribution is interpretability and verifiability of KG reasoning,
not raw accuracy.
\end{abstract}
